% =====================================================================
%  Week 1 progress report — MammAlps behaviour-recognition baseline
%  Standalone; compile with: tectonic --synctex week1.tex
% =====================================================================
\documentclass[11pt,a4paper]{article}

\usepackage[a4paper,margin=2.4cm]{geometry}
\usepackage{iftex}
\ifPDFTeX
  \usepackage[T1]{fontenc}
  \usepackage[utf8]{inputenc}
  \IfFileExists{lmodern.sty}{\usepackage{lmodern}}{}
\else
  \usepackage{fontspec}
\fi
\usepackage{enumitem}
\usepackage{booktabs}
\usepackage[dvipsnames]{xcolor}
\usepackage[colorlinks=true,linkcolor=NavyBlue,urlcolor=NavyBlue]{hyperref}

\setlist{nosep,leftmargin=1.4em}
\newcommand{\done}{\textcolor{ForestGreen}{\textbf{done}}}
\newcommand{\next}{\textcolor{BurntOrange}{\textbf{next}}}

\title{\textbf{MammAlps Behaviour-Recognition Baseline}\\[2pt]
\large Week 1 Progress Report}
\author{Kostiantyn Boiar \\ \small MSc Dissertation, University of Glasgow}
\date{\today}

\begin{document}
\maketitle

\section{Objective and scope}
The dissertation asks whether an animal's foreground \emph{silhouette shape} improves
multimodal behaviour recognition. Week~1 targets the mandatory first step: a trustworthy,
reproducible \textbf{video-only} and \textbf{video+audio} baseline on the
\textbf{MammAlps Benchmark~I} dataset, scored with the authors' official evaluation script.
The shape modality is added later; this week is about the foundation it bolts onto.

A hard design constraint shapes everything: the heavy encoders (VideoMAE for video, an
AudioMAE-style model for audio) are \textbf{frozen}. We never train them; we run them once,
cache their outputs to disk, and train only a small head on top. This is cheap and, crucially,
keeps later comparisons fair (the silhouette stream is judged on information, not extra
parameters).

\section{Approach and strategy}
The work is organised as a sequence of small steps, each with a concrete artifact and a check
before moving on. Heavy, run-once feature extraction is designed for a \textbf{Kaggle free
GPU} (run once); the light, repeatable head training and evaluation run on an \textbf{M4
laptop} (PyTorch MPS). Week~1 covers the setup work plus the first two checks (split
verification and the evaluation-format check); no model has been trained yet.

\section{Environment setup}
\begin{itemize}
  \item \textbf{Repository \& version control.} Initialised git; created the
        \texttt{mammalps\_b1/} source tree (config, data, encoders, models, scripts, utils)
        with a \texttt{.gitignore} that keeps data, caches, weights and virtual environments
        out of git. The project is on GitHub as \texttt{mambr}.
  \item \textbf{Local environment.} A dedicated \texttt{mammalps-env} virtual environment on
        Python~3.11 with PyTorch (MPS available), \texttt{transformers}, audio/video decoding
        and data libraries; \texttt{ffmpeg} installed. Dependencies pinned in
        \texttt{requirements-local.txt} and frozen to a lock file.
  \item \textbf{Official scorer vendored.} Copied \texttt{eval\_b1.py} and
        \texttt{labels\_mapping\_b1.json} from the MammAlps repository at a pinned commit, so
        scoring is reproducible. A wrapper script (\texttt{run\_eval.sh}) invokes it.
  \item \textbf{Kaggle templates.} A GPU setup-check cell and notes on staging the 88\,GB
        dataset, ready for the heavy extraction step.
  \item \textbf{Reproducibility scaffolding.} A seeding utility and a data-free
        \textbf{smoke test} that builds the fusion head on random tensors and confirms the
        toolchain end to end. \done
\end{itemize}

\section{The evaluation contract and dataset facts}
Before any modelling, the exact input/output contract of \texttt{eval\_b1.py} was read from
source and locked:
\begin{itemize}
  \item Predictions are a tab-separated file, one row per (test clip $\times$ sampled
        segment), with all prediction columns before all label columns, each a bracketed
        float vector.
  \item Three tasks: \textbf{Species} (5 classes) and \textbf{Activity} (11) are multiclass;
        \textbf{Action} (19) is \textbf{multi-label}. This corrected an assumption in the
        original brief and determines the head's loss (sigmoid/BCE for actions, softmax/CE for
        the rest).
  \item The scorer groups rows by clip, averages the segments, and \textbf{asserts exactly
        1244 test clips}. This assertion is our format check.
\end{itemize}

\section{Data and splits}
\begin{itemize}
  \item \textbf{Data access without the 88\,GB download.} The per-clip label/split CSVs live
        \emph{inside} \texttt{mammalps\_v1.zip}. Since the host supports HTTP range requests,
        a small fetcher pulls only those files ($\sim$660\,KB) instead of the whole archive.
  \item \textbf{Splits verified.} Parsed the comma-separated CSVs
        (\texttt{video\_path, start\_s, end\_s, activity, actions, species}) into typed clip
        records with one-hot/multi-hot label vectors, and confirmed the counts:
        \textbf{train 4205 / val 686 / test 1244 / total 6135} --- matching the paper exactly.
  \item \textbf{Findings.} The data is heavily imbalanced (the test set is $\sim$98\%
        red deer; activities dominated by foraging and vigilance), and the contact/posture
        behaviours the project ultimately targets are rare (courtship~56, chasing~3,
        escaping~1). This justifies macro-mAP and class-balanced sampling, and is important
        context for the later per-behaviour analysis. \done
\end{itemize}

\section{Code architecture}
The codebase is \textbf{object-oriented} and \textbf{fully config-driven}: every constant
lives in a central \texttt{Config} (paths, task definitions, dataset facts, model dimensions,
training hyper-parameters) with defaults overridable via YAML or environment variables.
Feature code is class-based (\texttt{LabelSpace}, \texttt{SplitLoader},
\texttt{MultiModalHead}/\texttt{MultiTaskLoss}, and script classes such as
\texttt{SplitChecker} and \texttt{MetadataFetcher}); only genuine utilities remain plain
functions. Style is enforced with full type hints, Google-style docstrings, and \texttt{ruff}
formatting/linting. Unrelated exploratory files from an earlier dataset (CalMS21) were removed
so the repository is focused on MammAlps. \done

\section{Evaluation-format check}
A \texttt{ResultsEmitter} class writes a \emph{dummy} results file for all 1244 test clips
($\times$10 segments): \textbf{real labels, random predictions}, in the exact column order and
vector format the scorer expects. Running \texttt{eval\_b1.py} on it confirmed the file passes
the \textbf{1244-clip assertion} and produces metrics for all three tasks. The average mAP came
out near-random ($\approx$\,\textbf{0.09}), exactly as expected for random guesses --- proving
the entire input/output pipeline is correct before any GPU time is spent. The emitter doubles
as the skeleton for the inference script later. \done

\section{Frozen encoders: one-clip probe}
The first step that touches the real video and audio. Two frozen encoder wrappers were
written: a \texttt{VideoEncoder} (decodes 16 frames from a clip's mp4 and runs a frozen
VideoMAE) and an \texttt{AudioEncoder} (loads the clip's wav and runs a frozen AST), each
producing one mean-pooled embedding and never trained. A \texttt{ClipProbe} pulls a clip's
mp4 and wav out of the 88\,GB archive with the same HTTP-range trick used for the CSVs (each
$\sim$0.2\,MB), then runs both encoders. On the M4 laptop (MPS) it produced finite
\textbf{768-dimensional} video and audio embeddings for real clips. These dimensions
($D_v = D_a = 768$) are exactly what the feature cache and the fusion head will consume. A
useful find: each clip ships with a \emph{separate} audio wav, so no audio has to be
demultiplexed from the video. One small fix was needed --- the latest \texttt{torchaudio}
routes \texttt{load} through a backend that was not installed, so wav loading was switched to
\texttt{soundfile}. \done

\section{Reporting infrastructure}
In parallel, the writing toolchain was set up: the Tectonic LaTeX compiler was installed and
wired into the editor for live, build-on-save PDF previews. A compile-blocking bug in the
proposal (an unclosed \texttt{\textbackslash author}) was fixed, the affiliation corrected to
the University of Glasgow throughout, and a full Glasgow-style dissertation skeleton drafted
(title page, front matter, chapter structure, a written Background \& Related Work chapter, and
a bibliography). The report directory was then split into \texttt{proposal/},
\texttt{dissertation/} and \texttt{progress/} sections. \done

\section{Status and next steps}
\begin{center}
\begin{tabular}{@{}ll@{}}
\toprule
Item & Status \\
\midrule
Environment, vendored scorer, smoke test & \done \\
Splits parsed + verified (test=1244, total=6135) & \done \\
Evaluation-format check (scorer accepts our output) & \done \\
One-clip frozen-encoder probe (768-d video + audio) & \done \\
Full feature cache on Kaggle (run once) & \next \\
Train head: video, then video+audio (target: V+A $>$ V) & --- \\
Multi-seed report; lock the baseline & --- \\
\bottomrule
\end{tabular}
\end{center}

\paragraph{Immediately next.} Run the same frozen encoders over all 6135 clips on a Kaggle
GPU --- sampling about ten windows per clip --- and write the cached embeddings (one file per
clip) plus a manifest, ready for the light head to train on.

\paragraph{Open risks / decisions.}
\begin{itemize}
  \item \textbf{Frozen vs fine-tuned.} The paper fine-tunes its encoder; our frozen-feature
        probe may score below its reported $\sim$0.45. The acceptance bar this step is a clean
        run plus \emph{video+audio $>$ video}; matching the exact number is a stretch goal.
  \item \textbf{Encoders.} Using nearest-available pretrained checkpoints rather than the
        paper's exact weights; absolute numbers may shift.
  \item \textbf{Kaggle storage.} The 88\,GB dataset exceeds Kaggle's working disk; a staging
        strategy must be settled before the full extraction (shard-on-the-fly, or a hosted
        dataset).
\end{itemize}

\end{document}
